\section[Pregunta 5: Atributos de calidad y trade-offs]{Pregunta 5: Comparación de atributos de calidad y trade-offs\\(sistema de transacciones financieras)}

% Responsable: Miguel Andrés Cazorla Zamora

El CTO plantea tres requisitos que no apuntan en la misma dirección: consistencia absoluta, agilidad de despliegue y viabilidad con un equipo de experiencia mixta. La decisión arquitectónica consiste precisamente en identificar cuál de ellos es innegociable y cuáles admiten una solución parcial. En un sistema de transacciones financieras, el primero lo es: un asiento contable que quede a medias no es una inconsistencia temporal, es dinero mal contabilizado.

\subsection{Requisito 1: consistencia de datos absoluta (ACID)}

\paragraph{Service-Based Architecture.} Los servicios de dominio comparten una base de datos, de modo que existe un único límite transaccional y la atomicidad la garantiza el gestor de base de datos \parencite{richardsford2020}. Un adeudo y su correspondiente abono se escriben dentro de la misma transacción: o constan ambos o no consta ninguno, sin código adicional. Las cuatro propiedades ACID se obtienen del motor, no del diseño de la aplicación.

\paragraph{Microservices Architecture.} El estilo exige que cada servicio posea su propia base de datos; esa es la condición que hace posible el despliegue verdaderamente independiente \parencite{richardsford2020}. En consecuencia, una operación que afecte a dos servicios no puede ser atómica. La alternativa pasa por una \emph{saga}: transacciones locales encadenadas con compensaciones, bajo un modelo BASE de consistencia eventual \parencite{fordhardparts2021}. Esto significa que existe un intervalo ---por breve que sea--- en el que el adeudo está registrado y el abono no.

\paragraph{Veredicto.} \textbf{Service-Based Architecture cumple el requisito; Microservices Architecture no.} La consistencia eventual no es una versión más lenta de la consistencia fuerte: es una garantía distinta, que admite estados intermedios visibles. Se puede construir un sistema financiero correcto sobre sagas, pero entonces el requisito deja de ser «consistencia absoluta» y pasa a ser «consistencia eventual con compensaciones auditadas», que es otra cosa. La objeción habitual ---hacer un microservicio lo bastante grande como para contener la transacción--- es en realidad un argumento a favor de SBA: un servicio de grano grueso con base de datos compartida es, por definición, el estilo que estamos comparando.

\subsection{Requisito 2: despliegue independiente}

Aquí conviene deshacer un malentendido frecuente: SBA no es un monolito. Sus servicios de dominio se despliegan de forma independiente, y esa es exactamente la propiedad que lo distingue de una arquitectura monolítica en capas \parencite{richardsford2020}. Un cambio en el servicio de informes no obliga a desplegar el de liquidación.

El matiz está en la base de datos compartida, que introduce un punto de acoplamiento estático: un cambio de esquema puede afectar a varios servicios y forzar un despliegue coordinado. La independencia es, por tanto, real pero condicionada ---independiente «con asterisco»---, y el asterisco se activa solo en los cambios de esquema, no en los cambios de lógica, que son la mayoría.

En microservicios la independencia es completa, porque cada servicio posee sus datos y nada obliga a coordinar. El coste está en otra parte: ese grado de autonomía requiere canalizaciones de integración continua maduras, automatización de pruebas de contrato, versionado de interfaces, observabilidad distribuida y una cultura operativa que lo sostenga \parencite{richardsford2020}. La independencia de despliegue no es gratuita; se paga en infraestructura y en disciplina de equipo.

El acoplamiento del esquema, además, es mitigable con técnicas conocidas: asignar la propiedad de cada conjunto de tablas a un único servicio ---los demás acceden a través de él---, usar esquemas lógicos separados dentro de la misma base de datos (el caso de base de datos «parcialmente compartida» que contempla el propio estilo) y aplicar migraciones en dos fases de tipo \emph{expand/contract}, que permiten desplegar los servicios en cualquier orden.

\subsection{Requisito 3: equipo de experiencia mixta}

Este requisito, que a primera vista parece el más blando, es el que en la práctica decide, y refuerza la misma conclusión.

Microservices Architecture es el estilo con menor puntuación en simplicidad y mayor coste de todos los analizados por \textcite{richardsford2020}. Exige dominar contenedores, orquestación, trazas distribuidas, patrones de resiliencia y, sobre todo, sagas. Y una saga mal implementada en un sistema financiero no produce un fallo visible: produce una compensación que no se ejecuta, un mensaje duplicado que contabiliza dos veces o un estado intermedio que alguien consulta. Son errores silenciosos, difíciles de reproducir y con consecuencias contables. Pedir a un equipo de experiencia mixta que escriba correctamente compensaciones distribuidas en el dominio financiero es, sencillamente, un riesgo mal calibrado.

SBA, en cambio, es de los estilos distribuidos el más próximo a lo que un equipo con experiencia desigual ya sabe hacer: pocos servicios de grano grueso, transacciones de base de datos convencionales, despliegue por servicio. Mantiene la modularidad y la testabilidad sin exigir la madurez operativa de los microservicios.

\subsection{Síntesis y trade-offs que el CTO acepta}

\begin{table}[h]
\centering
\small
\begin{tabularx}{\textwidth}{L{3.4cm}XX}
\toprule
\textbf{Requisito} & \textbf{Service-Based} & \textbf{Microservices} \\
\midrule
1. Consistencia ACID & Garantizada por la base de datos compartida & No alcanzable; sagas y consistencia eventual \\
2. Despliegue independiente & Sí, condicionado por el esquema compartido & Sí, completo, con alto coste de infraestructura \\
3. Equipo de experiencia mixta & Complejidad operativa moderada & Complejidad y riesgo elevados \\
\midrule
\textbf{Resultado} & \textbf{Cumple los tres} & Cumple el segundo, incumple el primero \\
\bottomrule
\end{tabularx}
\caption{Evaluación de ambos estilos frente a los tres requisitos del CTO.}
\end{table}

La recomendación es \textbf{Service-Based Architecture}, y los trade-offs que el CTO asume al priorizar consistencia sobre agilidad son los siguientes:

\begin{itemize}[leftmargin=1.5em]
  \item \textbf{Coordinación en los cambios de esquema.} Las migraciones que afecten a tablas compartidas requerirán planificación conjunta. Es un coste recurrente y acotado, mitigable con \emph{expand/contract}.
  \item \textbf{Escalado de grano grueso.} No se podrá escalar una funcionalidad concreta de forma aislada, sino el servicio completo que la contiene; y la base de datos compartida fija un techo de escalabilidad que los microservicios no tienen.
  \item \textbf{Radio de impacto mayor.} Un defecto en un servicio de grano grueso afecta a más funcionalidad que uno en un microservicio, y la base de datos compartida constituye un punto único de fallo. El sistema forma un solo \emph{quantum} arquitectónico, de modo que todos los servicios comparten las mismas características operativas y no es posible dar a uno de ellos una disponibilidad o una elasticidad distintas \parencite{richardsford2020}.
  \item \textbf{Menor agilidad en la evolución.} Entre dos cambios equivalentes, el ciclo será algo más lento que en una arquitectura de microservicios bien operada.
\end{itemize}

Conviene cerrar señalando que la decisión no es irreversible. No todo el sistema tiene el mismo requisito de consistencia: la notificación al cliente, la generación de informes o la detección de fraude toleran consistencia eventual y pueden extraerse más adelante como servicios con sus propios datos, dejando el núcleo transaccional en SBA. Esta evolución selectiva es lo que \textcite{fordevolutionary2017} describen como arquitectura evolutiva, y recomiendan sostenerla con \emph{fitness functions} automatizadas ---por ejemplo, una verificación en la canalización de integración que impida que un servicio acceda a tablas de las que no es propietario--- para que el acoplamiento por la base de datos compartida no se degrade con el tiempo hasta convertir el sistema en el monolito que se quería evitar.
